\section{Experiments}
\label{sec:experiments}
The experiments test whether refinement preserves tracking, which
initialisation changes help, and how map quality relates to cost.
\subsection{Evaluation protocol}
We evaluate rendering on all eight Replica scenes~\cite{straub2019replica}
and TUM \texttt{freiburg1\_desk}~\cite{sturm2012tum}, and trajectory accuracy
on all nine TUM freiburg1 sequences. All compared maps and trajectories were
produced by local runs. Photo-SLAM and MonoGS use monocular input.
MASt3R-SLAM and VGGT-SLAM supply trajectory comparisons only.
We did not obtain a monocular MonoGS Replica map with the available
configurations; its depth-assisted Replica configuration is not included.
Our head-adaptation and refinement ablations additionally cover
EuRoC~\cite{burri2016euroc}, ETH3D~\cite{schops2017eth3d}, and
7-Scenes~\cite{shotton2013scenecoord}.

\paragraph{Common rendering conditions}
For each sequence we select 100 frames outside the union of participating
systems' keyframes and render each exported map using the same implementation.
The images are exposure-normalised and evaluated at $512{\times}288$
on Replica and $512{\times}384$ on TUM.
All exported SH bands are retained: degree three for Photo-SLAM and degree
zero for the evaluated MonoGS and our maps.
We align estimated camera centres to ground truth with
Umeyama's $\mathrm{Sim}(3)$ fit~\cite{umeyama1991}.
If the map-to-GT alignment is $(s,R,t)$, the camera-to-world rendering pose in
map coordinates is
\begin{equation}
  R_{\rm map}=R^\top R_{\rm gt},\qquad
  t_{\rm map}=R^\top(t_{\rm gt}-t)/s.
  \label{eq:eval}
\end{equation}
We report scene PSNR as $-10\log_{10}$ of the mean frame MSE and average
LPIPS-AlexNet~\cite{zhang2018lpips} across frames. Replica means weight each
scene equally. ATE is translation RMSE after similarity alignment.

\paragraph{Evaluation scope}
We measure \emph{non-keyframe reconstruction}. Tracked non-keyframes may
supervise refinement, so scoring views are not a strict unseen-view split.
The aligned GT cameras expose map and trajectory errors together.
Input preprocessing and optimisation history follow each system.

\paragraph{Configuration and budget}
External comparisons use the released Splatt3R head, fixed Gaussian centres,
no insertion thinning, and 300\,s post-sequence refinement.
Photo-SLAM runs to its own shutdown schedule; its optimisation budget is not
matched on Replica. MonoGS uses a 339\,s colour-refinement phase on TUM desk,
close to our post-sequence budget, but total pipeline costs differ.
Experiments ran on a host with two RTX~A6000 GPUs.
Photo-SLAM was compiled with OpenCV~4.10 in place of its tested 4.7/4.8 stack
for CUDA compatibility; no mapping algorithm was changed.
We retain per-frame metrics, input-artifact hashes, alignment transforms, and
selected frame indices for the new rendering evaluation.

\paragraph{Component controls}
Paired ablations retain the evaluation settings of each recorded study;
they are separate from the fresh full-SH external comparison.
The TUM, EuRoC, and 7-Scenes heads were trained for 40 epochs,
with batch sizes 2, 4, and 8 and learning rates $10^{-5}$,
$1.41{\times}10^{-5}$, and $2{\times}10^{-5}$. The 15\% validation split
is within each dataset family. Refinement is active and thinning disabled
in both arms of the head comparison.

\subsection{Reconstruction quality}
\begin{table}[t]
  \centering
  \caption{Replica reconstruction, monocular input, full exported SH.
  Bold marks the better value within each scene and metric.}
  \label{tab:replica}
  \small\setlength{\tabcolsep}{4pt}
  \begin{tabular}{@{}lrrrr@{}}
\toprule
& \multicolumn{2}{c}{PSNR (dB) \up} & \multicolumn{2}{c}{LPIPS \dn}\\
Scene & Photo-SLAM & Ours & Photo-SLAM & Ours\\
\midrule
office0 & 22.19 & \textbf{26.30} & 0.219 & \textbf{0.104} \\
office1 & 17.77 & \textbf{22.08} & 0.190 & \textbf{0.115} \\
office2 & 20.15 & \textbf{20.97} & \textbf{0.135} & 0.150 \\
office3 & 19.20 & \textbf{20.18} & 0.165 & \textbf{0.144} \\
office4 & 16.43 & \textbf{23.65} & \textbf{0.129} & 0.156 \\
room0 & 17.81 & \textbf{25.44} & 0.159 & \textbf{0.110} \\
room1 & 21.68 & \textbf{21.91} & 0.181 & \textbf{0.139} \\
room2 & 20.64 & \textbf{23.62} & \textbf{0.123} & 0.160 \\
\midrule
Mean & 19.48 & \textbf{23.02} & 0.163 & \textbf{0.135} \\
\bottomrule
\end{tabular}

\end{table}

Table~\ref{tab:replica} reports 23.02\,dB mean PSNR for our maps and
19.48\,dB for Photo-SLAM, a 3.54\,dB difference. Our scene PSNR is higher
on all eight sequences, whereas LPIPS improves on five of eight
(0.135 versus 0.163 on average).
The PSNR margin varies from 0.23\,dB on room1 to 7.64\,dB on room0.
Photo-SLAM has better LPIPS on office2, office4, and room2, showing that
the PSNR ordering does not imply uniformly better perceptual reconstruction.
Fig.~\ref{fig:evidence}(a) displays this scene dependence.

\begin{table}[t]
  \centering
  \caption{TUM desk reconstruction. Ours uses 300\,s polish;
  MonoGS uses 339\,s colour refinement. $N$ counts exported primitives.}
  \label{tab:tum}
  \small
  \begin{tabular}{@{}lrrr@{}}
\toprule
Method & PSNR (dB) \up & LPIPS \dn & $N$ (thousands)\\
\midrule
Photo-SLAM & 10.67 & 0.562 & 36.1 \\
MonoGS & 13.91 & 0.395 & 23.0 \\
Ours & 13.95 & 0.411 & 2396.9 \\
\bottomrule
\end{tabular}

\end{table}

On TUM desk (Table~\ref{tab:tum}), PSNR is close to MonoGS:
13.95 versus 13.91\,dB. MonoGS has lower LPIPS, 0.395 versus 0.411.
Both outperform Photo-SLAM on this sequence.

\begin{figure*}[t]
  \centering
  \includegraphics[width=\textwidth]{qualitative_tum.pdf}
  \caption{\textbf{Related-work comparison on TUM desk.}
  Each method occupies one column; each view is followed by an identically
  located crop (orange box). The two views are the middle ranks of the
  per-frame Ours-minus-MonoGS PSNR difference among the 100 common evaluation
  frames. All columns use the same target camera and image scale.
  The close quantitative comparison is reflected in remaining blur and
  geometric distortion in both MonoGS and our renderings.}
  \label{fig:qualitative}
\end{figure*}

\begin{figure*}[t]
  \centering
  \includegraphics[width=\textwidth]{evidence.pdf}
  \caption{\textbf{Quality distribution and map budget.}
  (a) Fresh full-SH scene PSNR differences on Replica; o/r denote office/room.
  (b) A separate recorded study retaining the highest-opacity primitives in
  four of our maps, without re-optimisation. The final marker is each full
  map, containing 2.0--3.0 million Gaussians. The study measures post-hoc
  removal from maps optimised at full size.}
  \label{fig:evidence}
\end{figure*}

Fig.~\ref{fig:qualitative} shows the two middle-ranked views by per-frame
PSNR difference against MonoGS. The teaser uses the same rule against
Photo-SLAM on office0. All methods share crop coordinates.

\subsection{Trajectory accuracy}
\begin{table}[t]
  \centering
  \caption{TUM freiburg1 ATE RMSE (m), similarity aligned.
  MASt3R denotes the inherited MASt3R-SLAM tracking control.
  Best values are determined over all five systems.}
  \label{tab:ate}
  \small\setlength{\tabcolsep}{3.2pt}
  \begin{tabular}{@{}lrrrrr@{}}
    \toprule
    Sequence & Ours & MASt3R & VGGT & Photo & MonoGS\\
    \midrule
    360   & 0.042 & 0.048 & 0.050 & \best{0.035} & 0.177\\
    desk  & 0.017 & 0.016 & 0.025 & \best{0.015} & 0.036\\
    desk2 & 0.028 & \best{0.024} & 0.029 & 0.438 & 0.844\\
    floor & 0.027 & 0.025 & 0.099 & \best{0.013} & 0.539\\
    plant & \best{0.015} & 0.020 & 0.025 & 0.046 & 0.071\\
    room  & \best{0.059} & 0.061 & 0.064 & 0.510 & 0.791\\
    rpy   & \best{0.022} & 0.023 & 0.026 & 0.057 & 0.041\\
    teddy & 0.048 & 0.045 & \best{0.036} & 0.305 & 0.123\\
    xyz   & \best{0.009} & \best{0.009} & 0.014 & 0.010 & 0.017\\
    \midrule
    Mean  & \best{0.030} & \best{0.030} & 0.041 & 0.161 & 0.293\\
    \bottomrule
  \end{tabular}
\end{table}
Table~\ref{tab:ate} shows the same rounded mean trajectory error as
MASt3R-SLAM, supporting the separation of geometric tracking and mapping.
Photo-SLAM has the lowest error on several sequences but larger errors on
desk2, room, and teddy; MonoGS also exhibits large errors on several runs.
The means are consequently sensitive to those sequences.
Multi-threaded tracking is not deterministic: two Photo-SLAM desk runs yielded
0.0174 and 0.0149\,m ATE. Single-run differences should be interpreted
with this variability in mind.

\subsection{Shared anchors and refinement}
\label{sec:ablations}
\paragraph{Pose-correction control}
A blockwise $\mathrm{Sim}(3)$ perturbation applies a 6\% scale correction
to saved TUM desk and room anchors. With fixed world-space supervision
cameras, refinement largely reverses that correction after about 500 steps.
When supervision follows the anchors, the correction is retained.
This tests (\ref{eq:consistency}) under a controlled update.

\begin{table}[t]
  \centering
  \caption{Paired component ablations from separate recorded runs.
  Negative relative $\Delta$LPIPS indicates improvement.}
  \label{tab:ablation}
  \small\setlength{\tabcolsep}{4pt}
  \begin{tabular}{@{}llrr@{}}
    \toprule
    Component & Sequence & $\Delta$PSNR & $\Delta$LPIPS\\
    \midrule
    \multicolumn{4}{@{}l}{\emph{Refiner on vs.\ off, released head, 300\,s polish}}\\
    Refinement & ETH3D sofa\_1 & $+8.69$ & $-58.3\%$\\
               & Replica office0 & $+4.62$ & $-56.8\%$\\
               & TUM desk & $+3.52$ & $-27.3\%$\\
               & EuRoC V1\_01 & $+1.00$ & $-14.5\%$\\
    \midrule
    \multicolumn{4}{@{}l}{\emph{Family head vs.\ released head, refinement in both}}\\
    Adaptation & EuRoC MH\_01 & $+0.62$ & $-20.8\%$\\
               & TUM desk & --- & $-7.9\%$\\
               & TUM room & --- & $-9.9\%$\\
               & 7-Scenes chess & $+0.10$ & $+0.14\%$\\
    \bottomrule
  \end{tabular}
\end{table}
\paragraph{Refinement and head adaptation}
The first block of Table~\ref{tab:ablation} holds the released head fixed
and toggles refinement, with gains of 1.00--8.69\,dB across four datasets.
These values include 300\,s polish and are not instantaneous online gains.
The second block holds refinement active and disables thinning in both arms,
isolating head adaptation. Improvement is substantial on EuRoC and TUM,
while 7-Scenes chess changes little and its LPIPS slightly worsens.
Thus family-specific adaptation is useful but not uniformly beneficial.

\subsection{Initialisation and online budget}
\paragraph{Insertion prior}
Table~\ref{tab:thinning} lists all 12 recorded thinning comparisons.
Ten cells improve LPIPS clearly and two are near zero.
The confidence-weighted versus uniform allocation has only a
$-0.06\%$ median LPIPS difference, with its interval spanning zero.
Opacity attenuation helps more consistently than the confidence allocation.
Without thinning, refinement also reduces the median opacity of a
released-head TUM map from about 1.0 to 0.26.

\begin{table}[t]
  \centering
  \caption{Insertion ablation: relative LPIPS change for
  $\lambda=0.45$ versus off. Refinement is active in both arms.}
  \label{tab:thinning}
  \small\setlength{\tabcolsep}{4pt}
  \begin{tabular}{@{}llr@{}}
    \toprule
    Scene & Head & $\Delta$LPIPS\\
    \midrule
    Replica office3 & Plain adapted & $-12.7\%$\\
    Replica office0 & Plain adapted & $-12.6\%$\\
    Replica office2 & Plain adapted & $-10.2\%$\\
    Replica room0 & Plain adapted & $-8.9\%$\\
    Replica room1 & Plain adapted & $-7.1\%$\\
    Replica office1 & Plain adapted & $-1.6\%$\\
    Replica office0 & Opacity knee 0.9 & $-5.3\%$\\
    Replica office0 & Opacity knee 0.6 & $-1.8\%$\\
    TUM desk & Released & $-6.4\%$\\
    TUM desk & Adapted & $-0.6\%$\\
    EuRoC MH\_01 & Adapted & $-5.6\%$\\
    7-Scenes chess & Adapted & $+0.16\%$\\
    \bottomrule
  \end{tabular}
\end{table}

\paragraph{Causal delivery}
In causal replay, images enter supervision only when they arrive.
Table~\ref{tab:budget} compares this schedule with refinement after mapping.
Most of the gain appears by 500 steps; 3000 steps give no further causal
improvement in this TUM desk control.
A separate dual-GPU run gains 2.15\,dB with tracker median latency changing
from 101 to 103\,ms; sharing one GPU raises it to 206\,ms.
These are focused measurements on TUM desk.
In a broader GUI-enabled, dual-GPU campaign using family heads, seven of nine
sequences improve LPIPS before polish. TUM room and 360 worsen LPIPS by
11.1\% and 20.6\%, despite small PSNR gains. Larger maps remain harder
to refine within the available online budget.

\begin{table}[t]
  \centering
  \caption{TUM desk iteration-budget control, PSNR (dB).
  Both arms use the recorded map-evaluation protocol.}
  \label{tab:budget}
  \small
  \begin{tabular}{@{}lrrrr@{}}
    \toprule
    Steps & 120 & 500 & 1000 & 3000\\
    \midrule
    Causal & 13.65 & 14.39 & 14.47 & 14.41\\
    Post-hoc & 13.80 & 14.32 & 14.35 & 14.37\\
    \bottomrule
  \end{tabular}
\end{table}

\subsection{Cost and map size}
\begin{table}[t]
  \centering
  \caption{Recorded system cost on TUM desk, using the same 1\,Hz GPU
  memory sampler. Ours here has refinement disabled; these are not the
  costs of the polished maps in Table~\ref{tab:tum}.}
  \label{tab:cost}
  \small
  \begin{tabular}{@{}lrr@{}}
    \toprule
    System & Wall time (s) & Peak GPU memory (MiB)\\
    \midrule
    Ours, refiner off & 70 & 21\,035\\
    Photo-SLAM & 28 & 1\,286\\
    MonoGS & 555 & 2\,389\\
    \bottomrule
  \end{tabular}
\end{table}
Our TUM desk map contains 2.397 million primitives versus 36\,069 for
Photo-SLAM and 23\,019 for MonoGS. Replica maps contain 2.0--3.3 million.
Table~\ref{tab:cost} reports about $16.4\times$ Photo-SLAM's peak GPU memory
before adding our refinement cost. Photo-SLAM is also faster in this
end-to-end measurement. The broader GUI-enabled deployment campaign used up
to 41.1\,GB on the tracking GPU plus 14.0\,GB on the refinement GPU.

To measure how much the representation depends on its size, we retain the
highest-opacity Gaussians at decreasing budgets and re-render without further
optimisation. Fig.~\ref{fig:evidence}(b) shows that quality remains closer to
the full map around one million primitives but drops markedly below that.
For example, office0 falls from 26.29\,dB to 20.23\,dB at 300\,000
primitives in this separate study. This loss motivates optimising directly
at a target map size.
